\section{Threefold contractions for the boundary argument}\label{sec:finite-null}

The boundary argument uses the following specialization of
Das--Hacon--P\u{a}un's threefold contraction theorem
\cite[Theorem~5.5]{DasHaconPaun2024}.  It supplies contractions both on
the floors of the supported fourfold program and on lower strata in the
restriction argument.  Throughout these preparations, a real
\((1,1)\)-class on a singular space is a Bott--Chern class defined by
smooth local potentials.

\begin{proposition}[Threefold contraction of an adjoint plus a K\"ahler class]
\label{prop:threefold-contraction}
Let \((Z,\Gamma)\) be a normal connected compact K\"ahler klt threefold
with effective rational boundary and actual \(\Q\)-Cartier adjoint
\(A=K_Z+\Gamma\).  Suppose
\begin{equation}\label{eq:threefold-exposed-input}
 \alpha=c_1(A)+[\omega]\quad\hbox{is nef},\qquad
 \omega\quad\hbox{is K\"ahler}.
\end{equation}
Then there is a projective surjection \(f:Z\to Y\) with
\(f_*\OO_Z=\OO_Y\), where \(Y\) is normal compact K\"ahler with
rational singularities, and a K\"ahler class \([\omega_Y]\) satisfying
\(\alpha=f^*[\omega_Y]\).  The actual rational line \(-A\) is
\(f\)-ample.  A compact curve is contracted exactly when its
\(\alpha\)-degree is zero.
\end{proposition}

The proposition assumes no bigness of \(\alpha\), pseudo-effectivity
of \(A\), or factoriality of \(Z\).  Its proof is given in
Subsection~\ref{subsec:threefold-contraction-bridge}, after the required
constructions.  The separate exposure input
\cite[Corollary~5.3]{DasHaconPaun2024} supplies a class of the form
\eqref{eq:threefold-exposed-input} exposing any negative extremal ray in
the cited threefold cone theorem.  In that case the proposition contracts
exactly the curves in that ray.

The preparation has three parts.  First we contract the entire finite
curve null locus of a nef and big class, obtaining a normal compact
analytic target.  Next, conormal direct-image vanishings and analytic
thickenings extend a contraction of a prime floor to its ambient space.
Finally, these two constructions combine with the named projective MMP,
descent, positivity, and nonbig inputs to prove the proposition and its
floor restriction.  The K\"ahler target in the proposition is part of
this combined argument, beyond the finite-null construction alone.

\subsection{Contraction of the full finite null locus}

We first prove the finite-null contraction assertion of
\cite[Proposition~6.2]{DasHacon2024} used in the nef and big cases.
For an irreducible positive-dimensional analytic subspace \(V\), its top
intersection with such a class is computed on a resolution of \(V\).
Write
\[
 \operatorname{Null}(\alpha)=
 \bigcup_{\substack{V\subset X\ {\rm irreducible}\\
                     \dim V>0,\ \alpha^{\dim V}\cdot V=0}} V .
\]
\begin{proposition}[A finite null locus]\label{prop:finite-null}
Let \(X\) be a normal connected ordinary \(\Q\)-factorial compact
K\"ahler threefold with klt singularities.  Let \(\alpha\) be a nef and
big real \((1,1)\)-class.  Suppose that
\(C=\operatorname{Null}(\alpha)\) is a finite union of curves.  There is
a proper bimeromorphic morphism \(\phi:X\to X'\) to a normal compact
analytic space such that \(\phi\) is an isomorphism on \(X\setminus C\),
and the underlying sets of its nontrivial fibers are exactly the
connected components of \(C\).
\end{proposition}
\begin{proof}
The assertion is the identity map if \(C\) is empty.  Otherwise give
\(C\) its reduced structure and write \(C_1,\ldots,C_s\) for its
irreducible components.

We will construct a positive line on a resolution and use it to give
the conormal of the full analytic inverse image a positive presentation.
Grauert's contraction criterion will then apply to that conormal.

\subsubsection*{An exceptional divisor and one actual positive line}

Choose a projective resolution \(\pi:\widehat X\to X\) with smooth
source which is an isomorphism over \(X_{\rm reg}\) and has pure
divisorial exceptional locus
\cite[Theorem~2.13]{DasHaconPaun2024}.  Compactness makes the resolution
a finite composition of blowups.  Its source is compact K\"ahler:
metrics of a relatively ample line can be patched using a partition of
unity from the base, preserving positivity on the vertical tangent
spaces, and a sufficiently large pulled-back K\"ahler form makes the
curvature positive in every direction.  Compactness supplies one
constant for this last step.

Put \(\beta=\pi^*\alpha\).  Pullback of the metric nef approximations
shows that \(\beta\) is nef.  We check the positive top intersection
needed on the smooth source.  By bigness choose a K\"ahler current
\(T_X\in\alpha\) with \(T_X\geq\omega_X\), where \(\omega_X\) is a
K\"ahler form on \(X\).  Pulling back its local plurisubharmonic
potentials gives
\[
 T=\pi^*T_X\in\beta,\qquad T\geq\eta:=\pi^*\omega_X.
\]
The pulled potentials are not identically minus infinity, since
\(\widehat X\) dominates \(X\).  The form \(\eta\) is smooth and
semipositive.  Fix a K\"ahler form \(\kappa_{\widehat X}\), and for
\(\delta>0\) choose a positive form
\(b_\delta\in\beta+\delta[\kappa_{\widehat X}]\).  Wedge the positive
current \(T-\eta\), successively, with \(b_\delta^2\),
\(\eta\wedge b_\delta\), and \(\eta^2\).  Integration and
\(\delta\downarrow0\) give
\begin{equation}\label{eq:null-positive-cube}
 \beta^3\geq\beta^2[\eta]\geq\beta[\eta]^2\geq[\eta]^3>0 .
\end{equation}
The last inequality holds because \(\eta\) is positive on a nonempty
open set.  All these products are on the smooth compact \(\widehat X\);
only a current wedged with smooth semipositive forms was used.

For an irreducible \(d\)-dimensional positive-dimensional subspace
\(V\subset\widehat X\), projection of its fundamental cycle gives
\begin{equation}\label{eq:null-cycle-projection}
 \beta^d\cdot V=
 \begin{cases}
  0,&\dim\pi(V)<d,\\
  \deg(V/\pi(V))\,\alpha^d\cdot\pi(V),&\dim\pi(V)=d.
 \end{cases}
\end{equation}
The second case maps a null subspace into \(C\).  In the first case its
image is contained in the image of the exceptional locus, hence in
\(\operatorname{Sing}X\).  Normality of the threefold makes that singular
locus at most one-dimensional.  Thus every null subspace of \(\beta\)
maps into the set \(C\cup\operatorname{Sing}X\) of dimension at most one.

The smooth nef positive-volume criterion
\cite[Lemma~2.35]{DasHaconPaun2024} makes \(\beta\) big by
\eqref{eq:null-positive-cube}.  Collins--Tosatti
\cite[Theorem~1.1]{CollinsTosatti2015} identifies its non-K\"ahler locus
on \(\widehat X\) with \(\operatorname{Null}(\beta)\).  Choose the K\"ahler current
\(S\in\beta\) supplied by
\cite[Section~2.5.1 and Theorem~3.17(ii)]{Boucksom2004}.  In that
convention its analytic singularities are defined by one coherent ideal
\(\mathfrak a\) on \(\widehat X\) and one coefficient \(c>0\), with
\(E_+(S)=V(\mathfrak a)=\operatorname{Null}(\beta)\) as sets.
Principalize this global ideal by
a finite sequence \(g:Y\to\widehat X\) of blowups with smooth centers,
using \cite[Remark~2.14]{DasHaconPaun2024}.  The source \(Y\) is smooth
and compact K\"ahler.  If
\(\mathfrak a\OO_Y=\OO_Y(-D_{\mathfrak a})\), the logarithmic
presentation and Poincaré--Lelong, in the cited normalization, give
\[
 g^*S=\theta+[G],\qquad G=cD_{\mathfrak a}\geq0,
\]
where \(G\) is a finite real divisor supported over the singular set and
\(\theta\) is a global smooth closed form.
If \(S\geq\epsilon\kappa_{\widehat X}\), the equality off \(G\) and
continuity show that \(\theta\geq\epsilon g^*\kappa_{\widehat X}\)
everywhere.

Discarding blowups along Cartier centers, which are isomorphisms, choose
for this finite composition an effective
\(g\)-exceptional integral divisor \(F\) with \(\OO_Y(-F)\)
\(g\)-ample.  Such a divisor is obtained from the exceptional
tautological divisors, giving earlier pullbacks sufficiently large
positive weights.  A smooth representative \(v\) of \(c_1(F)\) can be
chosen so that \(g^*\kappa_{\widehat X}-\delta v\) is K\"ahler for some
\(\delta>0\): the curvature of \(-F\) is positive on the vertical
kernels, and compactness bounds the horizontal and mixed terms.  If
\(g\) is an isomorphism, take \(F=0\).  Consequently
\(\theta-\epsilon\delta v\) is K\"ahler.  With \(p=\pi g\), we obtain
\begin{equation}\label{eq:null-exceptional-decomposition}
 p^*\alpha=\kappa_0+[E],\qquad
 \kappa_0=[\theta-\epsilon\delta v]\ \text{K\"ahler},\qquad
 E=G+\epsilon\delta F\geq0 .
\end{equation}
Every prime of \(G\) maps into \(C\cup\operatorname{Sing}X\), and every
prime of \(F\) maps to a set of codimension at least two in
\(\widehat X\).  Hence every prime of \(E\) is \(p\)-exceptional.
The morphism \(p\) is projective: projective morphisms compose over the
compact base \(X\) by \cite[Remark~2.11]{DasHaconPaun2024}.

Openness of the K\"ahler cone on the smooth \(Y\) permits a nonnegative
rational divisor \(E'\) with the same exceptional prime support and
coefficients sufficiently close to those of \(E\) so that
\[
 \kappa'=p^*\alpha-[E']\quad\hbox{is K\"ahler}.
\]
Choose \(r>0\) clearing its denominators, and put
\begin{equation}\label{eq:null-positive-line}
 H=rE'\geq0,\qquad L=\OO_Y(-H).
\end{equation}
Thus \(H\) is an integral exceptional Cartier divisor and \(L\) is an
actual holomorphic invertible sheaf.

We verify relative ampleness on every full fiber of \(p\).  Let \(a\)
be a smooth representative of \(\alpha\) with local potentials on \(X\),
let \(\Theta_L\) be the curvature of a smooth metric on \(L\), and choose
a K\"ahler representative for \(\kappa'\).  Equality of Bott--Chern
classes on \(Y\) gives a global smooth function \(\varphi\) with
\[
 r\kappa'=r p^*a+\Theta_L+dd^c\varphi .
\]
Adjust the metric of \(L\) by \(\varphi\).  On a base chart
\(a=dd^c\rho\), adding \(r\rho\circ p\) to its local weights makes their
curvature \(r\kappa'\).  On a fiber, that added function is constant.
The restricted weights therefore have strictly plurisubharmonic ambient
extensions.  The same weights define positivity on the full possibly
nonreduced fiber: nilpotents do not change their values or the absolute
values of transition units.  The compact analytic positivity criterion
and the proper fiberwise criterion
\cite[Definition~2.8, Corollary~1.12, and Remark~3.2]{Fujino2026Fujiki}
now show that \(L\) is \(p\)-ample.

\subsubsection*{Numerical ampleness on the projective inverse image}

The compact reduced curve \(C\) is projective, even if it is reducible
or disconnected.  Indeed, choose one point on each irreducible component
which is smooth on all of \(C\) and is outside the other components.
Their sum is an effective Cartier divisor.  Its canonical section is
nonzero and has a zero on each positive-dimensional irreducible subspace
of \(C\).  The positivity lemma in Section~3.4 of
\cite{Grauert1962}, followed by Section~3.2, Satz~2, makes its line
positive and ample.  Fix a very ample line \(A\) on \(C\).

Put
\[
 D=(p^{-1}C)_{\red},\qquad f_D:D\to C,\qquad L_D=L|_D .
\]
The reduced base change of \(p\) is projective over \(C\).  Since \(C\)
is compact and projective and the final base is a point, the
compact-base projective composition assertion makes \(D\) projective
\cite[Remark~2.11]{DasHaconPaun2024}.  Chow and GAGA allow \(D\) and
its holomorphic invertible sheaves to be treated as a projective scheme
over \(\C\) \cite{Serre1956GAGA}.  The space \(D\) can be reducible,
nonnormal, and of mixed dimension.  The restriction \(L_D\) is the
restriction of the actual line \(L\); a component of \(D\) can lie
inside \(H\).

Choose an ample line \(A_D\) on \(D\) with a smooth curvature form
\(\Theta_{A_D}\).  For one integer \(q>0\) sufficiently large,
\[
 qr\kappa'|_D-\Theta_{A_D}
\]
is positive.  To see this on the singular space, extend the local
weights of \(A_D\) to finitely many relatively compact ambient charts
covering \(D\), and bound their Hessians by the ambient K\"ahler form.
This is the local-potential positivity convention.  If
\(\Gamma\subset D\) is an integral curve, its image is a point or one
of the \(C_l\); projection on the normalizations gives
\((p^*\alpha)\cdot\Gamma=0\).  Therefore
\begin{equation}\label{eq:null-numerical-margin}
 \deg_\Gamma(L_D^q\otimes A_D^{-1})
   =qr\kappa'\cdot\Gamma-\Theta_{A_D}\cdot\Gamma>0 .
\end{equation}
The class of \(L_D^q\otimes A_D^{-1}\) is consequently in the dual of
the closed cone of curves of the projective scheme \(D\).  The numerical
space is finite-dimensional.  Kleiman's line-bundle criterion
\cite[Section~4.8 and Theorem~4.10]{Fujino2009Fundamental} says that the
ample \(A_D\) is positive on every nonzero class of the closed curve
cone.  Its minimum on a compact unit slice is positive, so its class
is in the interior of the dual cone.
The sum of an element of the dual cone and an element of its interior
is again in the interior.  Thus \(L_D^q\), and hence \(L_D\), is ample.
The ample subtraction in \eqref{eq:null-numerical-margin} supplies the
uniform margin; strict positivity on individual curves alone would not
give this conclusion.

The ample line \(L_D\) is the positivity needed for the contraction.
The remaining ideal calculations transfer it to a positive presentation
of the conormal of the full analytic inverse image, including its
nonreduced structure.

\subsubsection*{The full scheme inverse image}

Let \(\mathfrak c=\mathcal I_C\), and define
\begin{equation}\label{eq:null-full-inverse}
 J=\mathfrak c\OO_Y=\im(p^*\mathfrak c\to\OO_Y),\qquad
 Z=V(J)=Y\times_X C,\qquad f:Z\to C.
\end{equation}
The space \(Z\) has its full analytic subspace structure, including
possible multiplicities and embedded components, and \(Z_{\red}=D\).
Put \(K_m=p_*L^m\) for \(m\geq1\).  It is a coherent ideal of
\(\OO_X\): effectivity of \(H\) gives \(L^m\subset\OO_Y\), and
\(p_*\OO_Y=\OO_X\) by normality and proper bimeromorphicity.

We use relative Serre vanishing in the following exact form.  For a
proper holomorphic map, a relatively ample invertible sheaf, a fixed
coherent twist, and a compact part of the base, its positive higher
direct images vanish there for every sufficiently large tensor power.
This is \cite[Definitions~2.1--2.2 and Proposition~2.5]{Ancona1982};
it permits nonreduced complex spaces and assumes they are separated
and countable at infinity.  Those conditions hold for the present spaces and the
relatively compact opens used next.

On a base open \(U\) choose generators \(g_1,\ldots,g_t\) of
\(\mathfrak c\).  Let \(Q_U\) be the coherent kernel in
\[
 0\longrightarrow Q_U\longrightarrow\OO_{p^{-1}U}^{\oplus t}
 \xrightarrow{(p^*g_1,\ldots,p^*g_t)}
 J|_{p^{-1}U}\longrightarrow0 .
\]
Choose finitely many such opens with compact subsets whose interiors
cover \(C\).  Relative Serre vanishing for
\(Q_U\otimes L^m\) on these compact subsets and for \(J\otimes L^m\)
over \(C\) gives one integer \(M\) for which the needed \(R^1p_*\)'s
vanish for every \(m\geq M\).  Tensor the displayed presentation by
\(L^m\) and push forward.  Inside \(K_m\), the image of \(K_m^{\oplus t}\)
is exactly multiplication by the \(g_i\).  Hence, at every stalk over
\(C\),
\begin{equation}\label{eq:null-ideal-product}
 p_*(J\otimes L^m)=\mathfrak cK_m\qquad(m\geq M).
\end{equation}
This is an image computation; it uses no projection formula for the
possibly nonflat ideal \(\mathfrak c\).

Push forward the exact sequence
\[
 0\longrightarrow J\otimes L^m\longrightarrow L^m
 \longrightarrow L^m|_Z\longrightarrow0 .
\]
The other \(R^1\) vanishing and \eqref{eq:null-ideal-product} give the
canonical isomorphism of coherent \(\OO_C\)-modules
\begin{equation}\label{eq:null-full-quotient}
 P_m:=K_m/\mathfrak cK_m
   =K_m\otimes_{\OO_X}\OO_C
   \xrightarrow{\ \sim\ } f_*(L^m|_Z),\qquad m\geq M .
\end{equation}
In particular the full inverse \(Z\), rather than only its reduction,
occurs in this equality.  No flatness or Cartier assumption on \(J\)
has been used.

\subsubsection*{A positive presentation of the conormal}

Let \(\mathfrak n=\ker(\OO_Z\to\OO_D)\).  This coherent nilradical
satisfies \(\mathfrak n^e=0\) for one \(e\), by local Noetherianity and
compactness of \(Z\).  For
\[
 T_m=L^m|_Z\otimes f^*A^{-2},
\]
the filtration \(\mathfrak n^jT_m\) has quotients
\[
 (\mathfrak n^j/\mathfrak n^{j+1})\otimes
 L_D^m\otimes f_D^*A^{-2},\qquad 0\leq j<e.
\]
Their coefficient sheaves are fixed and coherent on the projective
\(D\).  Serre vanishing for the ample \(L_D\) kills their \(H^1\) for
every sufficiently large \(m\).  Induction on this finite filtration
then gives
\[
 H^1(Z,T_m)=0\qquad(m\gg0).
\]
Only \(H^1\) of the subobject and quotient is needed at each step, so
the embedded nilpotent layers are included.

Write \(F_m=f_*(L^m|_Z)\).  Projection formula with the invertible
\(A^{-2}\) and the low-degree Leray sequence give
\begin{equation}\label{eq:null-base-vanishing}
 H^1(C,F_m\otimes A^{-2})
   =H^1(C,f_*T_m)\lhook\joinrel\longrightarrow H^1(Z,T_m)=0 .
\end{equation}
This does not assert vanishing of \(R^1f_*T_m\).  Embed
\(i:C\hookrightarrow\PP^n\) by \(A\).  By GAGA the coherent sheaf
\(F_m\) is algebraic.  The sheaf \(i_*(F_m\otimes A^{-1})\) is
\(0\)-regular: its degree-one condition is
\eqref{eq:null-base-vanishing}, and all higher conditions vanish since
its support has dimension at most one.  Castelnuovo--Mumford regularity
\cite[Tag~08A8, Lemma~33.35.12]{StacksProject} gives a surjection
\begin{equation}\label{eq:null-positive-presentation}
 A^{\oplus n_m}\longrightarrow F_m=P_m\longrightarrow0
 \qquad(m\gg0).
\end{equation}
Coherent torsion is allowed in this argument.

Fix one such \(m\).  We first check that \(C\subset V(K_m)\).
For each \(C_l\), the projective \(D\) contains an integral curve
\(\Gamma\) dominating \(C_l\): take a component dominating \(C_l\)
and cut by sufficiently general ample hyperplanes, taking the component
itself when it is a curve.  From
\((p^*\alpha)\cdot\Gamma=0\) and \(\kappa'\cdot\Gamma>0\) we obtain
\(E'\cdot\Gamma<0\).  If \(\Gamma\) were not contained in the support
of the effective Cartier divisor \(H\), its canonical section on the
normalization of \(\Gamma\) would give nonnegative degree, a contradiction.
Thus an irreducible component of \(H\) contains \(\Gamma\).  Its image
contains \(C_l\) and has dimension at most one by exceptionality, so
that image equals \(C_l\).  At every
point of \(C_l\) a unit germ on \(X\) pulls back to a unit and cannot
vanish along that component.  Hence \(K_m\) is a proper ideal at every
point of \(C\).

Outside \(p(H)\), the ideal \(K_m\) equals \(\OO_X\).  The image of each
exceptional prime has codimension at least two in the normal threefold.
It follows that \(V(K_m)_{\red}\) is a finite union of curves and points.
Since it contains every \(C_l\), each \(C_l\) is one of its irreducible
curve components.
Let \(B\) be the union of its irreducible components other than the
selected curves \(C_l\), and put \(\mathfrak b=\mathcal I_B\).  Define
\[
 I=(K_m:\mathfrak b^\infty).
\]
This is a coherent ideal with one global finite saturation exponent.
Indeed, for each integer \(j\geq0\), the ideal
\((K_m:\mathfrak b^j)\) is the kernel of the coherent
morphism
\(\OO_X\to\mathcal Hom(\mathfrak b^j,\OO_X/K_m)\).
The ascending chain stabilizes at each stalk by Noetherianity.  Equality
of two successive coherent ideals then holds on a neighborhood, and the
colon recursion gives all later equalities there.  Compactness supplies
one exponent for all of \(X\).

On \(X\setminus B\) we have \(I=K_m\).  On \(X\setminus C\), the analytic
Nullstellensatz puts a local power of \(\mathfrak b\) inside \(K_m\),
so \(I=\OO_X\).  Closure of the dense parts \(C_l\setminus B\) now gives
\[
 V(I)_{\red}=C,\qquad \sqrt I=\mathfrak c .
\]
The set \(B\cap C\) is finite.  The induced map
\[
 P_m=K_m/\mathfrak cK_m\longrightarrow Q:=I/\mathfrak cI
\]
is therefore an isomorphism away from finitely many points of \(C\).
After tensoring by \(A^{-2}\), its kernel and cokernel are still
point-supported.  Splitting through its image, the two long exact
sequences, \eqref{eq:null-base-vanishing}, and the vanishing of positive
cohomology of point-supported sheaves give
\(H^1(C,Q\otimes A^{-2})=0\).  Here a coherent sheaf on the projective
curve has no \(H^2\).  The same regularity argument as above gives
\begin{equation}\label{eq:null-conormal-presentation}
 A^{\oplus N_Q}\longrightarrow Q\longrightarrow0 .
\end{equation}
Since \(I\subset\mathfrak c\), right exactness of restriction gives the
actual conormal identity
\begin{equation}\label{eq:null-conormal}
 Q=I/\mathfrak cI=(I/I^2)\otimes_{\OO_X}\OO_C .
\end{equation}
Its support is all of \(C\).  At each point there, \(I\) is a nonzero
proper finitely generated ideal and \(\mathfrak c\) lies in the maximal ideal, so
Nakayama gives \(I/\mathfrak cI\ne0\).  Nonzeroness of the ideal also
follows from its one-dimensional zero set in the normal threefold.

\subsubsection*{The coherent conormal contraction criterion}

For a coherent ideal, Grauert's associated normal linear space is
defined by its local linear relations.  Restricting a finite presentation
of \(I\) to \(C\) presents the module in \eqref{eq:null-conormal}.
Thus the reduced associated space in Sections~3.6--3.7 of
\cite{Grauert1962} is
\[
 N_I=\bigl(\operatorname{Specan}_C\operatorname{Sym}Q\bigr)_{\red}.
\]
This definition does not require \(Q\) to be locally free.  The
surjection \eqref{eq:null-conormal-presentation} embeds \(N_I\) as a
closed reduced linear subspace of
\(\operatorname{Tot}((A^{-1})^{\oplus N_Q})\).

Put \(V_A=H^0(C,A)^*\).  The very ample embedding
\(C\hookrightarrow\PP(V_A)\), where \(\PP(V_A)\) parametrizes lines,
identifies \(A^{-1}\) with the restricted
tautological line.  Consequently
\(\operatorname{Tot}((A^{-1})^{\oplus N_Q})\) is closed in
\(C\times V_A^{N_Q}\).  Its projection to \(V_A^{N_Q}\) is proper because \(C\)
is compact, and it is an analytic embedding off the zero section:
on the open set where its \(j\)-th vector is nonzero, projectivization
to that vector's line, followed by the inverse of the embedding of
\(C\) on its image, recovers the base point holomorphically.  The
squared Euclidean norm of this projection
is strictly plurisubharmonic off the zero section.  Its positive
sublevels are relatively compact neighborhoods of that section with
strongly pseudoconvex boundaries; scalar multiplication gives a nonzero
radial derivative at every positive level.  Restricting them to \(N_I\) proves
exactly the weak negativity of this reduced linear space, including
when \(C\) is singular or disconnected and \(Q\) has torsion.

Grauert's Section~3.7, Satz~8, applies to any coherent ideal with the
given compact zero set and with weakly negative associated normal
linear space \cite{Grauert1962}.  It does not require the ideal to be
radical or locally principal.  The ambient \(X\) is reduced and
normal, and its compact zero set \(C\) is a union of curves with no
isolated points.  The criterion therefore makes \(C\) exceptional.
Section~2, Definition~3 and Satz~5, of that source give a global proper
surjective holomorphic map \(\phi:X\to X'\) with \(\phi(C)\) discrete,
\(\phi^{-1}(\phi(C))=C\) as underlying sets, and
\[
 X\setminus C\simeq X'\setminus\phi(C).
\]
The Remmert reduction used there has connected fibers.  Hence each
connected component of \(C\) has one image point and two components
cannot have the same image.  Its normality statement for a normal source
(p.~337) makes the Remmert-reduced target neighborhood normal; gluing it to the
unchanged normal complement makes \(X'\) normal.  Compactness of \(X\)
makes \(X'\) compact, and the complement isomorphism makes \(\phi\)
bimeromorphic.  This proves the proposition.
\end{proof}

The proposition supplies the normal compact analytic contraction and
its fiber sets.  The later assertions that a target is K\"ahler and that
a class descends to a K\"ahler class are separate inputs.

\begin{lemma}[Finiteness when there is no null surface]\label{lem:no-null-surface}
Let \(X\) be a normal connected compact K\"ahler threefold, and let
\(\alpha\) be nef and big.  If no irreducible surface has zero top
intersection with \(\alpha\), then \(\operatorname{Null}(\alpha)\)
is a finite union of curves, possibly empty.
\end{lemma}
\begin{proof}
Use a projective resolution \(\pi:\widehat X\to X\) with smooth compact
K\"ahler source as above.  The class \(\beta=\pi^*\alpha\) is nef with
positive cube by \eqref{eq:null-positive-cube}.  Collins--Tosatti makes
\(\operatorname{Null}(\beta)\) a proper analytic set.

Each irreducible component \(Z\) of this analytic set is itself null.
It has positive dimension, since each of its points lies on a
positive-dimensional null subspace.
Suppose instead that \(\beta^{\dim Z}\cdot Z>0\).  Resolve \(Z\)
projectively.  The restricted pulled class on the smooth compact
K\"ahler resolution is nef with positive top intersection, so
Collins--Tosatti gives it a proper null set.  Choose a point of \(Z\)
outside the image of that set, the target nonisomorphism locus of the resolution,
the singular locus of \(Z\), and the other components of
\(\operatorname{Null}(\beta)\).  Such a point exists because these are
proper analytic subsets of \(Z\).  By the definition of
\(\operatorname{Null}(\beta)\), a null irreducible subspace passes through
it.  That subspace lies in \(Z\), since the point is on no other
component.  Its strict transform is null for the restricted class,
contradicting the choice of the point.

There are finitely many components of the compact analytic null set.
By \eqref{eq:null-cycle-projection}, each of its surface components is
exceptional, since a nonexceptional one would map to a null surface on
\(X\).  Thus their images, as well as those of its curve components,
have dimension at most one.  For any downstairs null curve \(C_0\),
choose an irreducible component \(V\) of \(\pi^{-1}C_0\) dominating
\(C_0\).  Proper surjectivity supplies such a component.  If
\(\dim V=1\), cycle projection gives
\(\beta\cdot V=\deg(V/C_0)\,\alpha\cdot C_0=0\); if \(\dim V>1\),
it gives zero by dimension drop.  Thus \(V\) lies in an upstairs null
component whose image contains \(C_0\).  Every downstairs null curve
is therefore contained in the finite union of images.  An irreducible
curve in a finite union of curves and
points is one of the curve components.  There are therefore only
finitely many downstairs null curves, as asserted.
\end{proof}
